\part{Das Problem}

\chapter{Was Lernen heißt}

\section{Der Leitfall}

\begin{merksatz}
\boxtitel{Sepsis-Frühwarnung aus einem einzigen Wert}
Auf der Intensivstation des Klinikums wird bei jedem Patienten das Laktat bestimmt. Etwa jeder zehnte entwickelt in den folgenden Stunden eine Sepsis.

\begin{center}
{\small
\renewcommand{\arraystretch}{1.5}
\begin{tabular}{@{}P{2.4cm}P{2.2cm}P{3.2cm}@{}}
\toprule
\textbf{Gruppe} & \textbf{Anteil} & \textbf{Laktat in mmol/l}\\
\midrule
ohne Sepsis & $90\,\%$ & um $1{,}2$, Streuung $0{,}5$\\
mit Sepsis & $10\,\%$ & um $3{,}0$, Streuung $1{,}0$\\
\bottomrule
\end{tabular}}
\end{center}

\textbf{Die Frage:} Ab welchem Laktatwert soll gewarnt werden?

\begin{center}
\begin{tikzpicture}[scale=1.0]
\draw[grau!60, line width=0.7pt] (0,0) -- (7.2,0);
\draw[tief, line width=1.3pt, smooth, samples=80, domain=0:7]
  plot ({\x}, {2.6*0.9*exp(-((\x/1.4-1.2)^2)/(2*0.25))/(0.5*2.5066)*0.55});
\draw[akzent, line width=1.3pt, smooth, samples=80, domain=0:7]
  plot ({\x}, {2.6*0.1*exp(-((\x/1.4-3.0)^2)/(2*1.0))/(1.0*2.5066)*0.55*4});
\draw[grau!60, line width=0.8pt, dashed] (3.41,0) -- (3.41,2.0);
\node[font=\scriptsize, color=tief] at (1.1,1.9) {ohne};
\node[font=\scriptsize, color=akzent] at (5.0,0.75) {mit Sepsis};
\node[font=\scriptsize, color=grau, anchor=north] at (3.41,-0.1) {$2{,}44$};
\node[font=\scriptsize, color=grau, anchor=north] at (0,-0.1) {$0$};
\node[font=\scriptsize, color=grau, anchor=north] at (7.0,-0.1) {$5$};
\node[font=\scriptsize\itshape, color=grau] at (3.5,-0.75) {Laktat in mmol/l (Sepsiskurve vierfach überhöht)};
\end{tikzpicture}
\end{center}

\begin{center}
\begin{minipage}{0.88\textwidth}
\kursiv{Die beiden Verteilungen überlappen. \textbf{Es gibt keine Schwelle, die alle richtig einordnet} --- jede Wahl erkauft weniger Fehlalarme mit mehr übersehenen Fällen oder umgekehrt. Die Frage ist also nicht, wie man Fehler vermeidet, sondern welche man in Kauf nimmt.}
\end{minipage}
\end{center}
\end{merksatz}

\section{Die Bestandteile}

\begin{definition}
\boxtitelm{Definition 1.1 --- die Lernaufgabe}
\begin{center}
{\small
\renewcommand{\arraystretch}{1.55}
\begin{tabular}{@{}P{2.6cm}P{6.0cm}@{}}
\toprule
\textbf{Merkmale} & $x \in \mathcal X$ --- was man beobachtet, hier der Laktatwert.\\
\textbf{Zielgröße} & $y \in \mathcal Y$ --- was man vorhersagen will, hier Sepsis ja oder nein.\\
\textbf{Verteilung} & $P$ auf $\mathcal X \times \mathcal Y$ --- \textcolor{akzent}{unbekannt}.\\
\textbf{Regel} & $h: \mathcal X \to \mathcal Y$ --- auch \Begriff{Hypothese} oder \Begriff{Klassifikator}.\\
\textbf{Verlust} & $L(h(x),y)$ --- was ein Fehler kostet.\\
\textbf{Daten} & $(x_1,y_1),\ldots,(x_n,y_n)$ unabhängig aus $P$ gezogen.\\
\bottomrule
\end{tabular}}
\end{center}
\end{definition}

\begin{merksatz}
\boxtitel{Was an dieser Definition wichtig ist}
\begin{center}
\begin{minipage}{0.9\textwidth}
\textbf{Die Verteilung $P$ ist unbekannt.} Man sieht nur die $n$ Beispiele.
\end{minipage}
\end{center}

\begin{center}
\begin{minipage}{0.88\textwidth}
\kursiv{Das unterscheidet das Lernen von der Statistik aus Band 17. Dort nahm man an, $P$ gehöre zu einer bestimmten Familie --- etwa den Normalverteilungen --- und schätzte deren Parameter. \textbf{Hier nimmt man über $P$ möglichst wenig an}, und gerade deshalb ist die Frage schwer. Im Leitfall kennen wir $P$ nur, weil wir es vorgegeben haben; das Klinikum kennt es nicht.}
\end{minipage}
\end{center}

\schritt{Die letzte Zeile ist die stärkste Annahme}
\begin{center}
\begin{minipage}{0.88\textwidth}
\kursiv{\enquote{Unabhängig aus $P$ gezogen} heißt: Die Daten, aus denen man lernt, und die, auf die man das Gelernte anwendet, stammen aus \kursiv{derselben} Verteilung. \textbf{Das ist genau die Annahme, die beim Sepsis-System von 2021 verletzt war} --- und sie steht in fast jedem Satz dieses Bandes als stillschweigende Voraussetzung.}
\end{minipage}
\end{center}
\end{merksatz}

\begin{definition}
\boxtitelm{Definition 1.2 --- das Risiko}
Das \Begriff{Risiko} einer Regel ist ihr erwarteter Verlust:
\[ \Risk(h) = \E\bigl[L(h(X),Y)\bigr] . \]
Beim \Begriff{Null-Eins-Verlust} --- ein Fehler kostet eins, ein Treffer null --- ist das
\[ \Risk(h) = \Pr\bigl[h(X) \ne Y\bigr] , \]
die Wahrscheinlichkeit einer Fehleinordnung.
\end{definition}

\begin{merksatz}
\boxtitel{Das Ziel}
\begin{center}
\begin{minipage}{0.9\textwidth}
\textbf{Gesucht ist eine Regel mit kleinem Risiko --- auf Daten, die man noch nicht gesehen hat.}
\end{minipage}
\end{center}

\begin{center}
\begin{minipage}{0.88\textwidth}
\kursiv{Das ist der Kern des Titels. \textbf{Eine Regel, die die gesehenen Daten gut erklärt, ist leicht zu finden} --- man lernt sie einfach auswendig. Eine Regel, die auch auf ungesehenen Daten stimmt, ist das eigentliche Ziel, und die Differenz zwischen beiden ist der Gegenstand der Kapitel~2, 3, 7 und 8.}
\end{minipage}
\end{center}
\end{merksatz}

\section{Die bestmögliche Regel}

\begin{merksatz}
\boxtitel{Wenn man $P$ kennte}
Sei $\eta(x) = \Pr[Y=1 \mid X=x]$ die Wahrscheinlichkeit einer Sepsis beim Laktatwert $x$.

\begin{definition}
\boxtitelm{Definition 1.3 --- Bayes-Regel}
\[ h^{*}(x) = \begin{cases} 1 & \text{falls } \eta(x) \ge \tfrac12\\ 0 & \text{sonst} \end{cases} \]
\end{definition}

\schritt{Satz 1.4 --- sie ist die beste}
Für jede Regel $h$ gilt beim Null-Eins-Verlust
\[ \Risk(h^{*}) \ \le\ \Risk(h) . \]

\schritt{Beweis}
Man betrachtet einen festen Wert $x$. Die Wahrscheinlichkeit eines Fehlers bei gegebenem $x$ ist
\[ \Pr\bigl[h(x) \ne Y \mid X=x\bigr] = \begin{cases} \eta(x) & \text{falls } h(x)=0 \quad\text{(Sepsis übersehen)}\\ 1-\eta(x) & \text{falls } h(x)=1 \quad\text{(Fehlalarm)} \end{cases} \]

\kursiv{Das Kleinere der beiden ist} $\min\{\eta(x),\,1-\eta(x)\}$ --- und genau dieses wählt $h^{*}$: Ist $\eta(x) \ge \tfrac12$, so ist $1-\eta(x)$ das Kleinere, und $h^{*}$ sagt $1$.

\kursiv{Also gilt für jedes $x$:}
\[ \Pr\bigl[h^{*}(x) \ne Y \mid X=x\bigr] \ \le\ \Pr\bigl[h(x) \ne Y \mid X=x\bigr] \]

\kursiv{Mittelung über $x$} erhält die Ungleichung, denn ein Mittelwert kleinerer Zahlen ist kleiner. $\blacksquare$

\begin{center}
\begin{minipage}{0.9\textwidth}
\textbf{Man entscheidet sich an jeder Stelle für das Wahrscheinlichere --- und das ist global optimal.}
\end{minipage}
\end{center}
\end{merksatz}

\begin{merksatz}
\boxtitel{Der Leitfall gerechnet}
\schritt{Die Wahrscheinlichkeit einer Sepsis nach dem Messwert}
Nach dem Satz von Bayes (Band 4):
\[ \eta(x) = \frac{0{,}1\cdot f_1(x)}{0{,}1\cdot f_1(x) + 0{,}9\cdot f_0(x)} \]
mit den Dichten $f_0$ und $f_1$ der beiden Gruppen.

\begin{center}
{\small
\renewcommand{\arraystretch}{1.48}
\begin{tabular}{@{}P{2.2cm}P{2.2cm}P{3.2cm}@{}}
\toprule
\textbf{Laktat} & $\eta(x)$ & \\
\midrule
$1{,}0$ & $0{,}008$ & \\
$1{,}5$ & $0{,}021$ & \\
$2{,}0$ & $0{,}108$ & \\
$2{,}5$ & $0{,}590$ & \textcolor{akzent}{über der Hälfte}\\
$3{,}0$ & $0{,}973$ & \\
\bottomrule
\end{tabular}}
\end{center}

\schritt{Die Schwelle}
Die Bayes-Regel warnt, wo $\eta(x) = \tfrac12$ überschritten wird:
\[ x^{*} = 2{,}44 \ \text{mmol/l} \]

\schritt{Das Bayes-Risiko}
\[ \Risk^{*} = \Risk(h^{*}) = 3{,}47\,\% \]

\begin{center}
{\small
\renewcommand{\arraystretch}{1.5}
\begin{tabular}{@{}P{2.4cm}P{1.8cm}P{2.0cm}P{2.0cm}@{}}
\toprule
\textbf{Schwelle} & \textbf{Fehler} & \textbf{Empf.} & \textbf{Spezif.}\\
\midrule
$1{,}50$ & $25{,}4\,\%$ & $93{,}3\,\%$ & $72{,}6\,\%$\\
$2{,}00$ & $6{,}5\,\%$ & $84{,}1\,\%$ & $94{,}5\,\%$\\
$\mathbf{2{,}44}$ & $\mathbf{3{,}5\,\%}$ & $71{,}4\,\%$ & $99{,}3\,\%$\\
$3{,}00$ & $5{,}0\,\%$ & $50{,}0\,\%$ & $100\,\%$\\
\midrule
nie warnen & $10{,}0\,\%$ & $0\,\%$ & $100\,\%$\\
\bottomrule
\end{tabular}}
\end{center}

\begin{center}
\begin{minipage}{0.9\textwidth}
\textbf{Die beste mögliche Regel übersieht fast drei von zehn Sepsisfällen.}
\end{minipage}
\end{center}

\begin{center}
\begin{minipage}{0.88\textwidth}
\kursiv{Und sie ist nur um den Faktor drei besser als die Regel \enquote{nie warnen}, die ohne jede Messung auskommt. \textbf{Das Bayes-Risiko ist eine untere Schranke für jedes Verfahren} --- kein noch so kluges Lernverfahren kommt darunter, solange es nur das Laktat kennt. Wer bessere Vorhersagen will, braucht mehr Merkmale, nicht ein besseres Verfahren.}
\end{minipage}
\end{center}
\end{merksatz}

\section{Wenn Fehler verschieden kosten}

\begin{merksatz}
\boxtitel{Eine übersehene Sepsis ist schlimmer als ein Fehlalarm}
Kostet ein übersehener Fall $c_{\text{FN}}$ und ein Fehlalarm $c_{\text{FP}}$, so lautet die Bayes-Regel
\[ h^{*}(x) = 1 \quad\Longleftrightarrow\quad \eta(x) \ \ge\ \frac{c_{\text{FP}}}{c_{\text{FP}}+c_{\text{FN}}} . \]

\schritt{Warum}
Dieselbe Rechnung wie in Satz 1.4, nur mit Gewichten: Warnen kostet im Mittel $c_{\text{FP}}(1-\eta)$, nicht warnen $c_{\text{FN}}\,\eta$. Man warnt, wenn das Erste kleiner ist:
\[ c_{\text{FP}}(1-\eta) \le c_{\text{FN}}\,\eta \quad\Longleftrightarrow\quad \eta \ge \frac{c_{\text{FP}}}{c_{\text{FP}}+c_{\text{FN}}} \]

\schritt{Im Leitfall mit $c_{\text{FN}} : c_{\text{FP}} = 10 : 1$}
\[ \eta(x) \ \ge\ \frac{1}{11} = 0{,}091 \quad\Longrightarrow\quad x^{*} = 1{,}95 \ \text{mmol/l} \]

\begin{center}
{\small
\renewcommand{\arraystretch}{1.5}
\begin{tabular}{@{}P{2.6cm}P{1.8cm}P{1.8cm}P{2.0cm}@{}}
\toprule
& \textbf{Schwelle} & \textbf{Fehler} & \textbf{Kosten}\\
\midrule
kleinster Fehler & $2{,}44$ & $3{,}5\,\%$ & $0{,}292$\\
kleinste Kosten & $1{,}95$ & $7{,}4\,\%$ & $\mathbf{0{,}207}$\\
\bottomrule
\end{tabular}}
\end{center}

\begin{center}
\begin{minipage}{0.9\textwidth}
\textbf{Die Regel mit den geringsten Kosten macht doppelt so viele Fehler wie die mit den wenigsten Fehlern --- und ist dennoch die bessere.}
\end{minipage}
\end{center}

\begin{center}
\begin{minipage}{0.88\textwidth}
\kursiv{Die Empfindlichkeit steigt dabei von $71$ auf $85$ Prozent. \textbf{Die Wahl der Verlustfunktion ist keine technische Einzelheit, sondern eine medizinische Entscheidung} --- und sie wird in der Entwicklung solcher Systeme häufig nicht getroffen, sondern stillschweigend durch den Null-Eins-Verlust ersetzt. Das ist die erste Warnung des Vorworts in ihrer einfachsten Gestalt.}
\end{minipage}
\end{center}
\end{merksatz}

\section{Kein Lernen ohne Annahme}

\begin{merksatz}
\boxtitel{Satz 1.5 --- kein Verfahren ist überall besser}
Mittelt man über \kursiv{alle} möglichen Zusammenhänge zwischen $x$ und $y$, so haben alle Lernverfahren auf ungesehenen Daten dasselbe mittlere Risiko.

\schritt{Der Gedanke}
Sei $\mathcal X$ endlich und seien die Trainingsdaten gegeben. Auf einem Punkt $x$, der in den Daten nicht vorkommt, kann die Zielgröße jeden Wert haben.

\begin{center}
\begin{minipage}{0.9\textwidth}
\textbf{Ohne Annahme ist jede Vorhersage für $x$ genauso gut wie jede andere} --- im Mittel über alle denkbaren Welten liegt sie in der Hälfte der Fälle richtig.
\end{minipage}
\end{center}

\schritt{Ein kleines Beispiel}
Vier Laktatstufen, drei gesehen, eine nicht:
\begin{center}
{\small
\renewcommand{\arraystretch}{1.5}
\begin{tabular}{@{}P{2.2cm}P{1.4cm}P{1.4cm}P{1.4cm}P{1.4cm}@{}}
\toprule
\textbf{Stufe} & niedrig & mittel & erhöht & \textbf{hoch}\\
\midrule
gesehen & $0$ & $0$ & $1$ & \textcolor{akzent}{?}\\
\bottomrule
\end{tabular}}
\end{center}

\kursiv{Welt A:} \enquote{hoch} bedeutet Sepsis --- dann liegt eine Regel richtig, die \enquote{mehr Laktat, mehr Gefahr} unterstellt.\\
\kursiv{Welt B:} \enquote{hoch} bedeutet keine Sepsis --- dann liegt eine Regel richtig, die das Gegenteil annimmt.

\begin{center}
\begin{minipage}{0.88\textwidth}
\kursiv{Ohne Vorwissen sind beide Welten gleich möglich. \textbf{Das Vorwissen \enquote{mehr Laktat, mehr Gefahr} ist eine Annahme}, und erst sie macht die Vorhersage möglich. Jedes Lernverfahren trägt solche Annahmen in sich; man nennt sie seine \Begriff{induktive Neigung}. Sie ist kein Fehler, sondern die Voraussetzung dafür, überhaupt etwas zu lernen.}
\end{minipage}
\end{center}
\end{merksatz}

\begin{warnung}
\boxtitelr{Was der Satz nicht besagt}
\begin{center}
{\small
\renewcommand{\arraystretch}{1.52}
\begin{tabular}{@{}P{2.4cm}P{6.2cm}@{}}
\toprule
\textbf{Nicht:} & \enquote{Alle Verfahren sind gleich gut.} Auf den Aufgaben, die tatsächlich vorkommen, unterscheiden sie sich erheblich.\\
\textbf{Nicht:} & \enquote{Lernen ist unmöglich.} Es ist unmöglich \kursiv{ohne Annahme}.\\
\textbf{Sondern:} & Ein Verfahren ist gut, wenn seine Annahmen zur Aufgabe passen.\\
\bottomrule
\end{tabular}}
\end{center}

\begin{center}
\begin{minipage}{0.88\textwidth}
\kursiv{Das ist Humes Problem in mathematischer Gestalt --- und seine Lösung ist dieselbe wie im Alltag: \textbf{Man setzt voraus, dass die Welt nicht beliebig ist}, sondern Regelmäßigkeiten hat, die sich fortsetzen. Ob diese Voraussetzung im Einzelfall stimmt, lässt sich aus den Daten allein nicht entscheiden.}
\end{minipage}
\end{center}
\end{warnung}

\section{Der Bruchfall}

\begin{warnung}
\boxtitelr{Was der Leitfall verschweigt}
\begin{center}
{\small
\renewcommand{\arraystretch}{1.52}
\begin{tabular}{@{}P{2.4cm}P{6.2cm}@{}}
\toprule
\textbf{Die Verteilung} & Die Normalverteilungen waren vorgegeben. Tatsächliche Laktatwerte sind schief verteilt und nie negativ.\\
\textbf{Die Prävalenz} & Zehn Prozent gelten für diese Station. Auf einer Normalstation sind es vielleicht zwei.\\
\textbf{Die Zielgröße} & Was ist eine Sepsis? Die Definition hat sich 2016 geändert --- und damit, welche Fälle als $y=1$ zählen.\\
\bottomrule
\end{tabular}}
\end{center}

\schritt{Zur ersten Zeile}
Das Modell hat eine zweite Schwelle bei $x = -1{,}24$: Links davon würde die Bayes-Regel wieder warnen, weil die breitere Sepsisverteilung dort überwiegt.

\begin{center}
\begin{minipage}{0.88\textwidth}
\kursiv{Negative Laktatwerte gibt es nicht, und die Schwelle ist bedeutungslos. \textbf{Sie zeigt aber, dass ein Modell Aussagen über Bereiche macht, in denen es nie geprüft wurde} --- und dass diese Aussagen falsch sein können, ohne dass es auffällt, solange niemand dorthin schaut.}
\end{minipage}
\end{center}

\schritt{Zur zweiten Zeile}
Bei zwei Prozent Prävalenz statt zehn verschiebt sich die Bayes-Schwelle nach rechts, und die Fehlalarme nehmen relativ zu.

\begin{center}
\begin{minipage}{0.88\textwidth}
\kursiv{\textbf{Ein auf der Intensivstation gelerntes Verfahren ist auf der Normalstation ein anderes Verfahren}, auch wenn sich keine Zeile Programmcode ändert. Das ist das Problem des positiven Vorhersagewerts aus Band 17, Kapitel~5 --- hier als Problem der Übertragbarkeit.}
\end{minipage}
\end{center}

\schritt{Zur dritten Zeile}
\begin{center}
\begin{minipage}{0.88\textwidth}
\kursiv{Mit der Definition von 2016 änderte sich, welche Patienten als septisch gelten. Ein Verfahren, das auf Daten nach der alten Definition gelernt hat, sagt die alte Größe vorher. \textbf{Das Verfahren kann nicht klüger sein als die Etiketten, mit denen es gelernt hat} --- und Etiketten in der Medizin sind oft umstrittener, als ihre Darstellung in einer Tabelle vermuten lässt.}
\end{minipage}
\end{center}
\end{warnung}

\section{Übungsaufgaben}

\begin{aufgabe}
\boxtitel{Aufgabe 1.1 --- Die Bayes-Regel an Tabellenwerten}
Bei $1000$ Patienten wurde eine Laktatstufe erhoben:
\begin{center}
{\small
\renewcommand{\arraystretch}{1.5}
\begin{tabular}{@{}P{2.2cm}P{2.0cm}P{2.0cm}P{1.6cm}@{}}
\toprule
\textbf{Stufe} & \textbf{Sepsis} & \textbf{keine} & \textbf{Summe}\\
\midrule
niedrig & $5$ & $495$ & $500$\\
mittel & $15$ & $285$ & $300$\\
erhöht & $40$ & $110$ & $150$\\
hoch & $40$ & $10$ & $50$\\
\midrule
\textbf{Summe} & $100$ & $900$ & $1000$\\
\bottomrule
\end{tabular}}
\end{center}
Bestimmen Sie die Bayes-Regel und ihr Risiko.

\loesung
\schritt{Schritt 1 --- $\eta$ je Stufe}
\begin{center}
{\small
\renewcommand{\arraystretch}{1.48}
\begin{tabular}{@{}P{2.2cm}P{2.4cm}P{2.0cm}P{1.6cm}@{}}
\toprule
\textbf{Stufe} & \textbf{Rechnung} & $\eta$ & \textbf{Regel}\\
\midrule
niedrig & $5/500$ & $0{,}010$ & $0$\\
mittel & $15/300$ & $0{,}050$ & $0$\\
erhöht & $40/150$ & $0{,}267$ & $0$\\
hoch & $40/50$ & $0{,}800$ & $\mathbf{1}$\\
\bottomrule
\end{tabular}}
\end{center}

\begin{center}
\begin{minipage}{0.9\textwidth}
\textbf{Die Bayes-Regel warnt nur bei \enquote{hoch}.}
\end{minipage}
\end{center}

\schritt{Schritt 2 --- Die Fehler}
\begin{center}
{\small
\renewcommand{\arraystretch}{1.48}
\begin{tabular}{@{}P{2.2cm}P{2.4cm}P{2.4cm}@{}}
\toprule
\textbf{Stufe} & \textbf{Regel sagt} & \textbf{falsch}\\
\midrule
niedrig & keine & $5$\\
mittel & keine & $15$\\
erhöht & keine & $40$\\
hoch & Sepsis & $10$\\
\midrule
\textbf{Summe} & & $\mathbf{70}$\\
\bottomrule
\end{tabular}}
\end{center}

\[ \Risk(h^{*}) = \frac{70}{1000} = 7{,}0\,\% \]

\schritt{Schritt 3 --- Die Einordnung}
\begin{center}
{\small
\renewcommand{\arraystretch}{1.5}
\begin{tabular}{@{}P{2.6cm}P{1.6cm}P{2.0cm}P{1.6cm}@{}}
\toprule
\textbf{Regel} & \textbf{Fehler} & \textbf{übersehen} & \textbf{Fehlal.}\\
\midrule
nie warnen & $100$ & $100$ & $0$\\
Bayes (nur hoch) & $70$ & $60$ & $10$\\
ab erhöht & $135$ & $20$ & $120$\\
\bottomrule
\end{tabular}}
\end{center}

\kursiv{Zur dritten Zeile:} Übersehen werden $5+15 = 20$, Fehlalarme $110+10 = 120$.

\begin{center}
\begin{minipage}{0.9\textwidth}
\textbf{Die Bayes-Regel übersieht sechzig von hundert Fällen.} Ab \enquote{erhöht} zu warnen, halbiert die Übersehenen auf zwanzig --- verdoppelt aber die Gesamtfehlerzahl.
\end{minipage}
\end{center}

\schritt{Schritt 4 --- Mit Kosten $10:1$}
Die Schwelle liegt bei $\eta \ge \tfrac1{11} = 0{,}091$ --- nun wird ab \enquote{erhöht} gewarnt.
\begin{center}
\renewcommand{\arraystretch}{1.48}
\begin{tabular}{@{}P{2.6cm}P{3.0cm}P{1.8cm}@{}}
\toprule
\textbf{Regel} & \textbf{Kosten} & \\
\midrule
nur hoch & $10\cdot60 + 1\cdot10$ & $610$\\
ab erhöht & $10\cdot20 + 1\cdot120$ & $\mathbf{320}$\\
\bottomrule
\end{tabular}
\end{center}

\begin{center}
\begin{minipage}{0.88\textwidth}
\kursiv{Die Kosten halbieren sich, obwohl die Fehlerzahl sich fast verdoppelt. \textbf{Welche Regel die beste ist, entscheidet nicht die Mathematik, sondern die Bewertung der Fehler} --- und diese Bewertung gehört in die Hände derer, die die Folgen kennen.}
\end{minipage}
\end{center}
\end{aufgabe}

\begin{aufgabe}
\boxtitel{Aufgabe 1.2 --- Was die Prävalenz bewirkt}
Dieselben Laktatstufen werden auf einer Normalstation erhoben, wo nur zwei Prozent eine Sepsis entwickeln. Die Verteilung der Stufen \kursiv{innerhalb} jeder Gruppe bleibe gleich. Wie lautet nun die Bayes-Regel?

\loesung
\schritt{Schritt 1 --- Die bedingten Anteile aus Aufgabe 1.1}
\begin{center}
{\small
\renewcommand{\arraystretch}{1.48}
\begin{tabular}{@{}P{2.2cm}P{2.6cm}P{2.6cm}@{}}
\toprule
\textbf{Stufe} & $\Pr[\text{Stufe}\mid\text{Sepsis}]$ & $\Pr[\text{Stufe}\mid\text{keine}]$\\
\midrule
niedrig & $0{,}05$ & $0{,}550$\\
mittel & $0{,}15$ & $0{,}317$\\
erhöht & $0{,}40$ & $0{,}122$\\
hoch & $0{,}40$ & $0{,}011$\\
\bottomrule
\end{tabular}}
\end{center}

\kursiv{Zur ersten Zeile:} $5/100 = 0{,}05$ und $495/900 = 0{,}550$.

\schritt{Schritt 2 --- $\eta$ bei zwei Prozent}
\[ \eta = \frac{0{,}02\cdot p_1}{0{,}02\cdot p_1 + 0{,}98\cdot p_0} \]

\begin{center}
{\small
\renewcommand{\arraystretch}{1.48}
\begin{tabular}{@{}P{2.2cm}P{3.2cm}P{1.8cm}P{1.4cm}@{}}
\toprule
\textbf{Stufe} & \textbf{Rechnung} & $\eta$ & \textbf{Regel}\\
\midrule
niedrig & $\frac{0{,}001}{0{,}001+0{,}539}$ & $0{,}002$ & $0$\\
mittel & $\frac{0{,}003}{0{,}003+0{,}311}$ & $0{,}010$ & $0$\\
erhöht & $\frac{0{,}008}{0{,}008+0{,}120}$ & $0{,}063$ & $0$\\
hoch & $\frac{0{,}008}{0{,}008+0{,}011}$ & $\mathbf{0{,}424}$ & $\mathbf{0}$\\
\bottomrule
\end{tabular}}
\end{center}

\kursiv{Zur letzten Zeile:} $0{,}02\cdot0{,}40 = 0{,}008$ und $0{,}98\cdot0{,}01111 = 0{,}01089$; $\frac{0{,}008}{0{,}01889} = 0{,}424$.

\begin{center}
\begin{minipage}{0.9\textwidth}
\textbf{Die Bayes-Regel warnt nun gar nicht mehr.} Selbst bei hohem Laktat ist eine Sepsis weniger wahrscheinlich als keine.
\end{minipage}
\end{center}

\schritt{Schritt 3 --- Was das bedeutet}
\begin{center}
\begin{minipage}{0.88\textwidth}
\kursiv{Beim Null-Eins-Verlust ist \enquote{nie warnen} auf der Normalstation die beste Regel --- mit einem Risiko von zwei Prozent. \textbf{Das ist mathematisch korrekt und medizinisch unannehmbar}, und es zeigt, warum der Null-Eins-Verlust bei seltenen Ereignissen fast immer die falsche Wahl ist. Mit Kosten $10:1$ läge die Schwelle bei $0{,}091$, und \enquote{hoch} würde weiterhin warnen.}
\end{minipage}
\end{center}

\schritt{Schritt 4 --- Die allgemeine Lehre}
\begin{center}
\begin{minipage}{0.88\textwidth}
\kursiv{Die Messwerte innerhalb der Gruppen haben sich nicht geändert, die Regel schon vollständig. \textbf{Ein Verfahren, das auf einer Station mit anderer Prävalenz gelernt hat, trifft auf der neuen systematisch falsche Entscheidungen} --- und zwar in eine vorhersagbare Richtung. Das lässt sich korrigieren, wenn man die neue Prävalenz kennt; es bleibt unbemerkt, wenn man es nicht tut.}
\end{minipage}
\end{center}
\end{aufgabe}

\begin{aufgabe}
\boxtitel{Aufgabe 1.3 --- Auswendiglernen}
Ein Verfahren speichert alle $1000$ Patienten aus Aufgabe 1.1 samt Laborwert auf drei Nachkommastellen und sagt für jeden gespeicherten Wert das gespeicherte Ergebnis vorher. Für neue Werte sagt es \enquote{keine Sepsis}. Wie groß ist sein Fehler auf den Trainingsdaten, wie groß auf neuen Patienten?

\loesung
\schritt{Schritt 1 --- Auf den Trainingsdaten}
Jeder gespeicherte Patient wird richtig wiedergegeben --- sofern kein Wert zweimal mit verschiedenem Ergebnis vorkommt.

\[ \text{Trainingsfehler} \approx 0 \]

\schritt{Schritt 2 --- Auf neuen Patienten}
Ein neuer Laborwert stimmt auf drei Nachkommastellen fast nie mit einem gespeicherten überein. Das Verfahren sagt also fast immer \enquote{keine Sepsis}.

\[ \text{Testfehler} \approx 10\,\% \]

\begin{center}
\begin{minipage}{0.9\textwidth}
\textbf{Null Fehler auf den gesehenen Daten, zehn Prozent auf den ungesehenen --- genauso schlecht wie \enquote{nie warnen}.}
\end{minipage}
\end{center}

\schritt{Schritt 3 --- Der Vergleich}
\begin{center}
{\small
\renewcommand{\arraystretch}{1.5}
\begin{tabular}{@{}P{2.8cm}P{2.2cm}P{2.2cm}@{}}
\toprule
\textbf{Verfahren} & \textbf{Training} & \textbf{neu}\\
\midrule
Auswendiglernen & $\approx 0\,\%$ & $\approx 10\,\%$\\
Bayes-Regel (Stufen) & $7\,\%$ & $\approx 7\,\%$\\
\bottomrule
\end{tabular}}
\end{center}

\begin{center}
\begin{minipage}{0.88\textwidth}
\kursiv{Das ist das Grundproblem des Faches in seiner reinsten Form. \textbf{Das Verfahren mit dem kleinsten Trainingsfehler hat den größten Testfehler.} Es hat die Daten gelernt und nicht den Zusammenhang --- und genau das meint der Titel dieses Bandes: Man sucht nicht, was die gesehenen Daten erklärt, sondern was sich verallgemeinern lässt.}
\end{minipage}
\end{center}

\schritt{Schritt 4 --- Warum die Stufenregel verallgemeinert}
\begin{center}
\begin{minipage}{0.88\textwidth}
\kursiv{Sie fasst viele Patienten zu vier Gruppen zusammen und schätzt je Gruppe nur eine Zahl. Jede dieser Zahlen beruht auf mindestens fünfzig Beobachtungen. \textbf{Sie hat also wenig Freiheit} --- und wer wenig Freiheit hat, kann sich den Daten nicht zu sehr anpassen. Kapitel~2 macht aus diesem Gedanken eine Rechnung, Kapitel~7 einen Satz.}
\end{minipage}
\end{center}
\end{aufgabe}

\begin{geschichte}
\boxtitelg{Eine lernende Maschine und ein Name}
\Person{Alan Turing} schlug 1950 in seinem Aufsatz über das Denken von Maschinen vor, man solle nicht versuchen, einen erwachsenen Verstand nachzubauen, sondern den eines Kindes --- und ihn dann unterrichten.

\begin{center}
\begin{minipage}{0.9\textwidth}
\kursiv{Er nannte das eine \enquote{lernende Maschine} und sagte voraus, dass sie ihren Lehrer an manchen Stellen überraschen werde. \textbf{Die Vorhersage war richtig, und die Überraschungen sind bis heute nicht immer angenehm}: Lernende Maschinen finden in den Daten Muster, die ihre Entwickler nicht bemerkt hatten, und nicht alle davon sind die gewünschten.}
\end{minipage}
\end{center}

\schritt{Samuel}
\Person{Arthur Samuel} (1901--1990) schrieb ab 1952 bei IBM ein Dameprogramm, das durch Spiele gegen sich selbst besser wurde.

\begin{center}
\begin{minipage}{0.9\textwidth}
\kursiv{1959 veröffentlichte er einen Aufsatz, in dem der Ausdruck \enquote{maschinelles Lernen} erstmals erscheint. Er beschrieb es als Fähigkeit, zu lernen, ohne ausdrücklich programmiert zu sein. \textbf{Das Programm schlug 1962 einen starken Amateurspieler} --- ein Ereignis, das die Öffentlichkeit weit mehr beeindruckte als die Fachwelt.}
\end{minipage}
\end{center}

\schritt{Die Bayes-Regel ist älter}
Der Gedanke, an jeder Stelle für das Wahrscheinlichere zu entscheiden, stammt aus der statistischen Entscheidungstheorie.

\begin{center}
\begin{minipage}{0.9\textwidth}
\kursiv{\Person{Ronald Fisher} gab 1936 die lineare Diskriminanzanalyse an --- am Beispiel dreier Schwertlilienarten, deren Blütenmaße er vermessen ließ. Der Datensatz ist bis heute der meistbenutzte Übungsdatensatz des Faches. \textbf{Fisher war zugleich ein führender Eugeniker}, und die Zeitschrift, in der der Aufsatz erschien, trug diesen Begriff im Titel --- ein Umstand, der bei der Verwendung des Datensatzes lange übergangen wurde.}
\end{minipage}
\end{center}
\end{geschichte}

\begin{forschung}
\boxtitelp{Lernen in der Medizin heute}
\begin{center}
{\small
\renewcommand{\arraystretch}{1.45}
\begin{tabular}{@{}P{2.8cm}P{5.4cm}@{}}
\toprule
\textbf{Bildgebung} & Die erfolgreichste Anwendung --- Netzhautbilder, Hautläsionen, Röntgenbilder.\\
\textbf{Frühwarnung} & Sepsis, Verschlechterung, Wiederaufnahme --- gemischte Erfahrungen.\\
\textbf{Zulassung} & Lernende Medizinprodukte als eigene Regulierungsklasse.\\
\bottomrule
\end{tabular}}
\end{center}

\schritt{Zur ersten Zeile}
\begin{center}
\begin{minipage}{0.88\textwidth}
\kursiv{Bei Bildern erreichen Verfahren seit etwa 2017 in eng umrissenen Aufgaben die Treffsicherheit von Fachleuten. \textbf{Der Grund ist, dass die Zielgröße dort klar ist} --- eine Läsion ist bösartig oder nicht, und die Histologie entscheidet --- und dass die Bilder aus verschiedenen Kliniken einander ähnlicher sind als Laborverläufe.}
\end{minipage}
\end{center}

\schritt{Zur zweiten Zeile}
\begin{center}
\begin{minipage}{0.88\textwidth}
\kursiv{Frühwarnsysteme leiden unter allen drei Problemen des Bruchfalls: unklare Zielgröße, wechselnde Prävalenz, andere Datenlage an anderem Ort. \textbf{Die Nachprüfung von 2021 ergab eine Trennschärfe von $0{,}63$ statt der angegebenen $0{,}76$ bis $0{,}83$} --- und zwei Drittel der Sepsisfälle wurden nicht erkannt. Kapitel~3 behandelt, wie man eine solche Lücke vorher bemerkt.}
\end{minipage}
\end{center}
\end{forschung}

\begin{merksatz}
\boxtitel{Was Sie jetzt können}
\begin{center}
\renewcommand{\arraystretch}{1.5}
\begin{tabular}{@{}P{3.0cm}P{5.6cm}@{}}
\toprule
\textbf{Definition 1.1} & Merkmale, Ziel, unbekannte Verteilung\\
\textbf{Definition 1.2} & Risiko $=$ erwarteter Verlust\\
\textbf{Satz 1.4} & Bayes-Regel: überall das Wahrscheinlichere\\
\textbf{Kosten} & Schwelle $c_{\text{FP}}/(c_{\text{FP}}+c_{\text{FN}})$\\
\textbf{Satz 1.5} & kein Lernen ohne Annahme\\
\bottomrule
\end{tabular}
\end{center}

\begin{center}
\begin{minipage}{0.88\textwidth}
\kursiv{Die Bayes-Regel setzt voraus, dass man die Verteilung kennt. Das nächste Kapitel fragt, was geschieht, wenn man sie aus Daten schätzen muss --- und warum dabei zu viel Freiheit genauso schadet wie zu wenig.}
\end{minipage}
\end{center}
\end{merksatz}
